% The narration script of the explainer video, slide by slide: a handout for the presentation.
% script/body.tex and script/stats.tex are generated by video/tools/script_handout.py. Build from presentation/: make script
\documentclass[a4paper,11pt]{article}
\usepackage[a4paper,left=36mm,right=22mm,top=25mm,bottom=27mm,headsep=8mm,footskip=12mm,marginparwidth=26mm]{geometry}
\usepackage{amsmath}
\usepackage{fontspec}
\usepackage{unicode-math}
\setmainfont{Libertinus Serif}
\setmathfont{Libertinus Math}
\setmonofont{JetBrains Mono}[Scale=0.84]
\newfontfamily\calli{texgyrechorus-mediumitalic.otf}
\newfontfamily\garamond{EB Garamond}
\newfontfamily\stampfont{Libertinus Serif}[Numbers={Lining,Monospaced}]
\usepackage{microtype}
\usepackage{xcolor}
\usepackage{graphicx}
\usepackage{fancyhdr}
\usepackage{needspace}
\usepackage{url}
\usepackage{tikz}
\usepackage{tocloft}
\usepackage[hidelinks,bookmarksnumbered=false,pdfusetitle]{hyperref}

% the deck's palette, darkened for white paper
\definecolor{ink}{HTML}{1D1C24}
\definecolor{muted}{HTML}{6F6D79}
\definecolor{faint}{HTML}{C9C6CF}
\definecolor{gold}{HTML}{A2701A}
\definecolor{cyan}{HTML}{137B78}
\definecolor{rose}{HTML}{B3255B}
\color{ink}

\hypersetup{pdftitle={Uncommon Molecular Electronic Configurations: the narration script, slide by slide},
  pdfauthor={Shuvam Banerji Seal}, pdfsubject={CH5103 term paper: presentation handout},
  pdfkeywords={diradicals, Hund's rule, spin states, non-aufbau, narration script}}

\setlength{\parindent}{0pt}
\setlength{\parskip}{0pt}
\linespread{1.06}
\raggedbottom
\setlength{\emergencystretch}{2em}
\urlstyle{tt}
\setcounter{tocdepth}{2}

% ------------------------------------------------------------------ running heads
\pagestyle{fancy}
\fancyhf{}
\renewcommand{\headrulewidth}{0pt}
\fancyhead[L]{\garamond\scshape\footnotesize\color{muted}Narration script}
\fancyhead[R]{\garamond\footnotesize\color{muted}\leftmark}
\fancyfoot[C]{\stampfont\footnotesize\color{muted}\thepage}
\fancypagestyle{plain}{\fancyhf{}\fancyfoot[C]{\stampfont\footnotesize\color{muted}\thepage}}

% ------------------------------------------------------------------ contents
\renewcommand{\cfttoctitlefont}{\calli\Huge\color{gold}}
\renewcommand{\cftaftertoctitle}{\par\vspace{2mm}}
\renewcommand{\cftmarktoc}{\markboth{Contents}{}}
\renewcommand{\cftsecfont}{\garamond\scshape\color{gold}}
\renewcommand{\cftsecpagefont}{\stampfont\color{gold}}
\renewcommand{\cftsubsecfont}{\small}
\renewcommand{\cftsubsecpagefont}{\stampfont\small\color{muted}}
\setlength{\cftbeforesecskip}{9pt}
\setlength{\cftbeforesubsecskip}{0.5pt}
\setlength{\cftsubsecindent}{1.2em}
\setlength{\cftsubsecnumwidth}{2.4em}
\renewcommand{\cftsubsecdotsep}{\cftnodots}
\renewcommand{\cftsecdotsep}{\cftnodots}

% ------------------------------------------------------------------ the entries body.tex uses
\newcommand{\goldrule}[1]{\tikz\shade[left color=gold,right color=white] (0,0) rectangle (#1,0.6pt);}
\newcommand{\timespan}[2]{\stampfont\small\color{muted}#1\,–\,#2}

% \Part{Part I}{The paper}{start}{end}
\newcommand{\Part}[4]{%
  \clearpage
  \global\partopentrue
  \phantomsection
  \addcontentsline{toc}{section}{\texorpdfstring{#1\quad #2\quad{\stampfont\mdseries\color{muted}\small #3}}{#1: #2}}%
  \markboth{#1 · #2}{}%
  \vspace*{6mm}
  {\garamond\scshape\large\color{muted}#1\par}
  \vspace{2mm}
  {\calli\fontsize{40}{44}\selectfont\color{gold}#2\par}
  \vspace{3mm}
  \goldrule{9cm}\par
  \vspace{2mm}
  {\timespan{#3}{#4}\par}
  \vspace{9mm}}

% Notes-page layout: every slide starts a page (the first one of a part shares the part's opening page), the frame on
% top and its narration below. Each heading block is set in a box first, so the page break can be decided on its
% real height: the block and at least three lines of its script stay on one page. Video-only scenes run on.
\newif\ifpartopen
\newsavebox\scenebox
\newcommand{\placescene}{\par\needspace{\dimexpr\ht\scenebox+\dp\scenebox+3\baselineskip\relax}}

% \Slide{label: Slide n, Title slide or Closing slide}{n, or empty}{handout page}{title}{subtitle}{start}{end}
%        {caption: source frame and scene}{bookmark title}
\newcommand{\Slide}[9]{%
  \sbox\scenebox{\begin{minipage}{\textwidth}
    {\garamond\scshape\color{gold}#1}\hfill{\timespan{#6}{#7}}\par
    \vspace{0.6mm}
    {\large\bfseries #4\par}
    \if\relax\detokenize{#5}\relax\else\vspace{0.4mm}{\itshape\small\color{muted}#5\par}\fi
    \vspace{3mm}
    \centering
    \includegraphics[page=#3,width=0.8\textwidth]{ch5103-handout.pdf}\par
    \vspace{1.2mm}
    {\scriptsize\color{muted}#8\par}
  \end{minipage}}%
  \ifpartopen\global\partopenfalse\else\clearpage\fi
  \placescene
  \phantomsection
  \addcontentsline{toc}{subsection}{\protect\numberline{#2}\texorpdfstring{#4}{#9}}%
  \noindent\usebox\scenebox\par
  \vspace{4mm}}

% \VideoOnly{label}{title}{subtitle}{start}{end}{files it reads}{scene}{bookmark title}
\newcommand{\VideoOnly}[8]{%
  \sbox\scenebox{\begin{minipage}{\textwidth}
    {\garamond\scshape\color{cyan}#1\ \ ·\ \ video only}\hfill{\timespan{#4}{#5}}\par
    \vspace{0.6mm}
    {\large\bfseries #2\par}
    \if\relax\detokenize{#3}\relax\else\vspace{0.4mm}{\itshape\small\color{muted}#3\par}\fi
    \vspace{1.2mm}
    {\scriptsize\raggedright\color{muted}\if\relax\detokenize{#6}\relax\else reads #6\ \ ·\ \ \fi scene \texttt{#7}\par}
  \end{minipage}}%
  \placescene
  \phantomsection
  \addcontentsline{toc}{subsection}{\protect\numberline{#1}\texorpdfstring{#2}{#8}}%
  \noindent\usebox\scenebox\par
  \vspace{4mm}}

% \Line{time}{text}: one say() line, its start time in the left margin
\newcommand{\Line}[2]{%
  \noindent\llap{\stampfont\footnotesize\color{muted}#1\hspace{9pt}}#2\par\vspace{3.2pt}}

\newcommand{\EndScene}{\par\vspace{5mm}{\centering\color{faint}\rule{28mm}{0.4pt}\par}\vspace{6mm}}

\begin{document}
\input{script/stats}

% ------------------------------------------------------------------ title page
\pagenumbering{roman}
\begin{titlepage}
  \thispagestyle{empty}
  \vspace*{4mm}
  {\garamond\scshape\large\color{muted}CH5103 Term Paper · presentation handout\par}
  \vspace{10mm}
  {\calli\fontsize{40}{46}\selectfont\color{gold}Uncommon Molecular\\[1mm] Electronic Configurations\par}
  \vspace{5mm}
  {\Large Theory, derivations and verification of Gryn’ova, Coote \& Corminboeuf (2015)\par}
  \vspace{7mm}
  \goldrule{11cm}\par
  \vspace{7mm}
  {\calli\fontsize{26}{30}\selectfont The narration script, slide by slide\par}
  \vspace{3mm}
  {\large\itshape every line the explainer video speaks, word for word, beside the slide it narrates\par}
  \vspace{11mm}
  \includegraphics[page=1,width=0.72\textwidth]{ch5103-handout.pdf}\par
  \vfill
  {\large Shuvam Banerji Seal\par}
  \vspace{2mm}
  {\color{muted}\ScriptScenes\ scenes · \ScriptSlides\ slides · \ScriptLines\ narrated lines · \ScriptWords\ words ·
    {\stampfont\ScriptLength} of video\par}
  \vspace{1mm}
  {\small\color{muted}\url{https://github.com/Shuvam-Banerji-Seal/CH5103_Term_Paper}\par}
\end{titlepage}
\setcounter{page}{2}

% ------------------------------------------------------------------ how to read
\markboth{About this script}{}
{\calli\Huge\color{gold}About this script\par}
\vspace{6mm}

This is the complete narration of the explainer video \emph{When Hund’s Rule Fails} (release
\texttt{video-v1.1}, \ScriptLength), which walks through the CH5103 presentation on Gryn’ova, Coote and
Corminboeuf, “Theory and practice of uncommon molecular electronic configurations”, \emph{WIREs Comput.\ Mol.\ Sci.}\
\textbf{5}, 440–459 (2015), doi:10.1002/wcms.1233. Nothing has been edited or summarised: each paragraph is one line of
narration exactly as it was given to the voice, the Breeze-TTS-2 speech model, in the order it is spoken.

\vspace{4mm}
{\garamond\scshape\color{gold}Reading an entry}\par\vspace{1mm}
\begin{itemize}\setlength{\itemsep}{2pt}
  \item \textbf{Slide \emph{n}} is the frame number printed in the deck’s footer (\emph{n}\,/\,134); the title and
    closing slides carry no number. The picture is that frame from \texttt{ch5103-handout.pdf}, with every build
    step shown. Each of the deck’s 136 frames has its own scene in the video, in the same order.
  \item Each slide starts a new page, like a speaker’s notes page: the frame on top, its narration underneath.
  \item The time at the top right of an entry is when that scene starts and ends in the video. The time in the left
    margin is when the voice starts that line, rounded down to the second, so you can jump straight to it.
  \item The maths is written the way the voice says it (“M S equals zero”, “J a b minus K a b”). The slide above
    shows the same formulae typeset.
  \item Appendices A and B are in the video only (scenes A1–A18 and B1–B11). They read every verification script and every Lean proof line by
    line, so instead of a picture each entry lists the files it reads from the repository.
\end{itemize}

\vspace{3mm}
{\garamond\scshape\color{gold}How it was made, and checked}\par\vspace{1mm}
The document is generated by a script in the repository, \path{video/tools/script_handout.py}. It reads every
\texttt{say(…)} line from the scene sources in \texttt{video/scenes}, which is the exact text the voice was given. It puts the lines of each scene in spoken order
using the subtitle file written while that scene rendered, and finds each slide by counting frames in the deck’s
sources. The script refuses to write the document unless three checks pass. Every line must be found in its scene’s
subtitles, and those subtitles must contain nothing else. Every slide page must carry its frame’s title and footer number. The scene
lengths must add up to the video’s 322\,637 frames. Separately, all \ScriptLines\ margin times were compared with the
captions of the finished video, and every one agrees to within a second.

\vspace{3mm}
{\small\color{muted}The voice is synthetic: Breeze-TTS-2, used under its research, non-commercial licence.}

\clearpage
\markboth{Contents}{}
\tableofcontents

\clearpage
\pagenumbering{arabic}
\input{script/body}
\end{document}
